\documentclass[10pt,a4paper]{amsart}
\usepackage[margin=19mm]{geometry}
\usepackage{amsmath,amssymb,mathtools}
\usepackage[scr=rsfs,cal=euler]{mathalfa}
\usepackage{xfrac}
\usepackage[unicode,hidelinks,bookmarks=false]{hyperref}
\newcommand{\ie}{i.e.\ }
\newcommand{\cf}{cf.\ }
\newcommand{\resp}{respectively\ }
\newcommand{\Subsection}{\subsection}
\providecommand{\nobreakdash}{\nobreak\hbox{-}}
\setlength{\emergencystretch}{2em}
\multlinegap0pt
\usepackage{amssymb}
\newtheorem{lemma}{Lemma}
\newtheorem{thm}{Theorem}
\newtheorem{cor}{Corollary}
\newtheorem{Proof}{Proof}
\newcommand{\E}{\mathbf{E}}
\newcommand{\R}{\mathbb{R}}
\newcommand{\m}{\mathcal{M}}
\title{The Topology of Negatively Associated Distributions}
\author{Jonathan Root, Mark Kon}
\date{Source: arXiv:2304.09737v1 (CC BY 4.0)}
\begin{document}
\maketitle

\noindent\emph{Source and licence.} The Topology of Negatively Associated Distributions. Version arXiv:2304.09737v1 (CC BY 4.0), distributed by arXiv under the Creative Commons Attribution 4.0 International licence, http://creativecommons.org/licenses/by/4.0/, as recorded in the arXiv licence metadata for this exact version and shown on the abstract page https://arxiv.org/abs/2304.09737v1. Full licence notice: \texttt{licenses/arxiv-2304.09737-CC-BY-4.0.txt}.\par\smallskip

\noindent\textit{Editorial scope.} Counted unit: the article's cor cor (source0.tex lines 882--888). The article's complete original proof of it is delivered with it (source0.tex lines 890--1005, one \texttt{proof} environment of 116 lines). No mathematics is written, repaired or re-derived: the statement and the proof are the article's own lines, in the article's own notation, and no retained line is substituted or re-flowed.  Retained uncounted as the source-local context the proof needs: lines 245--252; lines 661--665; lines 852--854.  Nothing was omitted from any retained span, so the package carries no omission record.  No retained line contains a citation, so the wrapper carries no bibliography and no bibliography entry.  No retained line contains a figure environment, an image-inclusion command or drawing code, so no image is delivered and the package declares no asset dependency.  Editorial formatting only, in addition: an AMS \texttt{amsart} preamble that restates the article's own declarations and macros used by the retained lines, namely the environments \texttt{lemma}, \texttt{thm}, \texttt{cor}, \texttt{Proof} on separate counters, together with the standard packages those lines need.  The retained environments print in this document as Lemma 1, Theorem 1, Corollary 1, Proof 1 in order of appearance, whereas the article prints the same items with the numbering of its own full document.  The complete native source of the article as distributed in the arXiv e-print for this exact version is bundled unchanged as \texttt{sources/140318/source0.tex}.
\begin{lemma}\label{strictNAexist}
There exist strictly NA distributions on $\{0,1\}^n$
such that for all non-decreasing $f(x_I),g(x_J)$ (with $I,J$ disjoint sets of indices) having total variation 1 on $\{0,1\}^n$, there is an $\epsilon>0$ such that
\[
\int fg d\mu \leq \int f d\mu \int g d\mu  -\epsilon .
\]
Consequently, there exist strictly NA distributions on the whole of $\R^n$ satisfying the above equation (under the same measure $\mu$ supported on $\{0,1\}^n$ viewed as a subset of $\R^n$).
\end{lemma}

\begin{thm}\label{convexNA}
The space of negatively associated
distributions is not convex on $\R^n$. 
%However, the subset $E_{\mu,f}$, for any $\mu \in \m_{NA}(X^n)$, is convex.
\end{thm}

\begin{lemma}\label{NCmarginal}
For any $0<\epsilon<1$ and $1\leq i\leq n$, there exists a negatively correlated  distribution on the Boolean cube $I_n=\{0,1\}^n$ satisfying $\mu^{(i)}(1)=\epsilon$. In fact, given $0<\epsilon_i<1$, $i=1,\dots, n$, satisfying $\sum_i \epsilon_i=1$, there exist a negatively correlated distribution on the Boolean cube satisfying $\mu^{(i)}(1)= \epsilon_i$, $i=1,\dots, n$.
\end{lemma}

\begin{cor}
The space of negatively correlated distributions on the Boolean cube $I_n$ is non-convex. However, for any fixed $p_1,\dots, p_n \in \R$, the collection of measures
\[
E_{p_1,\dots, p_n} = \{\mu\in \m_{NC}(I_n) : \mu^{(i)}(1) = p_i, \; i=1,\dots, n\}
\]
is convex.
\end{cor}

\begin{Proof}
We begin with strictly negatively correlated distributions $\mu$ and 
$\nu$, which exist by Lemma \ref{strictNAexist}. We derive conditions under which the convex combination $\lambda\mu +(1-\lambda)\nu$ fails to be negatively correlated. We then produce strictly negatively correlated distributions whose convex combination fails to satisfy the above-mentioned condition.


Note that $\E_{\mu} x_i =\mu^{(i)}(1)$ and $\E_{\mu}x_ix_j =\mu^{(i,j)}(1,1)$ on the Boolean cube. Thus given strictly negatively correlated distributions $\mu$ and $\nu$, fix $i,j$ and note that 
\[
\mu^{(i,j)}(1,1) = \mu^{(i)}(1)\mu^{(j)}(1) -\epsilon_1
\]
for some $\epsilon_1>0$, and 
\[
\nu^{(i,j)}(1,1) = \nu^{(i)}(1)\nu^{(j)}(1) -\epsilon_2
\]
for some $\epsilon_2>0$.
Set $A= \mu^{(i)}(1)$, $B= \nu^{(i)}(1)$, $C= \mu^{(j)}(1)$,
and $D = \nu^{(j)}(1)$. Then if convexity is %*** Perhaps comment below should appear on p. 12
to hold, we once again must have (see proof of Theorem \ref{convexNA} above, where the roles of $f$ and $g$ are played by $x_i$ and $x_j$)
\[
\lambda AC +(1-\lambda)BD
-(\lambda\epsilon_1+(1-\lambda)\epsilon_2)
\leq \lambda^2 AC + \lambda(1-\lambda)AD
+\lambda(1-\lambda)BC + (1-\lambda)^2BD
\]
for $0\leq \lambda \leq 1.$ 
Now with $\tilde{C}=(A-B)(C-D)$ this simplifies to 
\[
\tilde{C}\lambda^2-\tilde{C}\lambda \geq -(\lambda \epsilon_1+(1-\lambda)\epsilon_2).
\]
As this holds for $\tilde{C}=0$, the second statement of the Corollary holds. 

Now set $\lambda=1/2$. We arrive at the condition 
\[
\frac{\tilde{C}}{4} \leq \frac{1}{2}(\epsilon_1+\epsilon_2).
\]
Our strategy is as follows. We introduce measures $\mu$ and $\nu$ as convex combinations of point masses at the standard basis vectors. We demonstrate that the inequality
\[
\frac{\tilde{C}}{4} \leq \frac{1}{2}(\epsilon_1+\epsilon_2)
\]
is valid. We then define a perturbation of this measure under which said inequality is violated, whereby we obtain a convex combination of measures which fail to be negatively correlated.

According to Lemma \ref{NCmarginal}, there exist strictly negatively correlated distributions with $\tilde{C}>0$. Specifically, define 
\[
\mu = \sum_k \beta_k \delta_{\alpha_k},
\]
and further
\[
\nu = \sum_k \beta_k' \delta_{\alpha_k}
\]
where $\alpha_k = (0,\dots, 1,\dots, 0)$ has exactly one 1 in the $k$th component, and $\sum _k \beta_k =1$, $\sum_k \beta_k'=1$. 
Then 
\[
\tilde{C} = (\beta_i -\beta_i')(\beta_j - \beta_j')
\]
where we have used that  $\tilde{C}=(A-B)(C-D)$ and $A= \mu^{(i)}(1)$, $B= \nu^{(i)}(1)$, $C= \mu^{(j)}(1)$,
and 
\[
\epsilon_1 = \beta_i\beta_j,
\;\;\;\epsilon_2 =\beta_i'\beta_j'.
\]
Thus the condition 
\[
\frac{\tilde{C}}{4} \leq \frac{1}{2}(\epsilon_1+\epsilon_2)
\]
becomes 
\[
\frac{(\beta_i -\beta_i')(\beta_j - \beta_j')}{4} \leq \frac{1}{2}(\beta_i\beta_j+\beta_i'\beta_j').
\]
%*** Try to show this inequality after \epsilon perturbation occurs
This inequality reduces to 
\[
\beta_i\beta_j +\beta_i'\beta_j' +\beta_i\beta_j' + \beta_i'\beta_j\geq 0,
\]
which holds for all non-negative reals.  %This proves the second statement of the Corollary regarding convexity of the subsets $E_{p_1,p_2,\ldots,p_n}.$

Since strict negative correlation is continuous in the Euclidean parameters of the distribution, we may perturb the above convex measure as follows. Define $\mu$ as above, but perturbed with a small additional weight given to $(1,1,0,\dots, 0)$: $\mu(1,1,0,\dots, 0) = \epsilon$. Here $\epsilon>0$ is small enough so that $\mu$ is still negatively correlated.  
Compensating this increased weight by a total decrease in the other decoupled positive weights totaling $\epsilon$ to keep normalization will still not affect negative correlation if $\epsilon$ is sufficiently small.
%***Do this in such a way that $\mu$ is still a probability distribution, and in such a way that $\tilde{C}$ does not change. Continued to 9/6/22 
Then $\mu^{(1,2)}(1,1) = \epsilon$ and $\mu^{(i)}(1) = \beta_i+\epsilon$ for $i=1,2$ . If we do the same for $\nu$, with the same $\epsilon$ perturbation, then note that $\tilde{C}$ does not change. Indeed, for $i=1, j=2$ 
\begin{eqnarray*}
    \tilde{C} &=& (A-B)(C-D) \\
    &=& (\beta_1+\epsilon-\beta_1'-\epsilon)(\beta_2+\epsilon -\beta_2'-\epsilon)\\
    &=&(\beta_1-\beta_1')(\beta_2-\beta_2').
    \end{eqnarray*}
%***Explain in a 5 word sentence why this is so
%***Also need to explain where we take away the \epsilon mass from the distribution so as to maintain measure 1, and to make sure we can make \epsilon as large as needed 
We may assume that $\tilde{C}>0$, by virtue of choosing $\beta_1>\beta_1'$ and $\beta_2>\beta_2'$.
On the other hand, for $\epsilon>0$ small enough
\[
 -\epsilon_1 \equiv \mu^{(1,2)}(1,1) -\mu^{(1)}(1)\mu^{(2)}(1) <0
\]
as the quantity 
\[
\mu^{(1,2)}(1,1) -\mu^{(1)}(1)\mu^{(2)}(1) =\epsilon-(\beta_1+\epsilon)(\beta_2+\epsilon)
\]
is continuous in $\epsilon$, and approaches $-\beta_1\beta_2$ from above as $\epsilon\to 0$. For $\beta_1$ and $\beta_2$ small enough, the quantity 
\[
\epsilon-(\beta_1+\epsilon)(\beta_2+\epsilon)
\]
will be positive; letting $\epsilon$ tend to 0 shows that this quantity will pass through zero, towards $-\beta_1\beta_2$. This shows that $\epsilon_1$ will move through 0, and thus the inequality 
\[
\frac{\tilde{C}}{4} \leq \frac{1}{2}(\epsilon_1+\epsilon_2)
\]
will be violated, as consideration of $\epsilon_2$ will be analogous. 
%consideration of the derivative in $\epsilon$ shows this to be the case: 
%the derivative in $\epsilon$ is $1-\beta_1-\beta_2\geq0$, using $\sum_k \beta_k=1$. 
%Thus $-\epsilon_1$ will increase towards 0 away from $-\beta_1\beta_2$ in $\epsilon$.



%***9/6/22 How big does \epsilon need to get to decrease this sufficiently, and will that still maintain negative correlation above
It follows that the right hand side of the inequality 
\[
\frac{\tilde{C}}{4} \leq \frac{1}{2}(\epsilon_1+\epsilon_2)
\]
will decrease. And thus we can, by continuity, decrease the right hand side until the inequality is violated. Thus it is shown that the space of negatively correlated distributions on the Boolean cube is non-convex.
\end{Proof}

\end{document}
